\documentclass[10pt]{beamer}

% --- Modern Theme ---
\usetheme{metropolis}           % A very modern, clean theme
\usecolortheme{default}
\setbeamercolor{background canvas}{bg=white}

% --- Packages ---
\usepackage[utf8]{inputenc}
\usepackage{booktabs}
\usepackage{array}
\usepackage{tabularx}
\usepackage[scale=2]{ccicons}
\usepackage{pgfplots}
\usepgfplotslibrary{dateplot}
\usepackage{xspace}

% --- Title Information ---
\title{Foundations of Artificial Intelligence}
\subtitle{Part I: Introduction \& Intelligent Agents}
\institute{Department of Computer Science}
\date{}

\begin{document}

\maketitle

\begin{frame}{Lecture Roadmap}
    \setbeamertemplate{section in toc}[sections numbered]
    \tableofcontents[hideallsubsections]
\end{frame}

\begin{frame}{Learning Objectives}
    By the end of this lecture, you should be able to:
    \begin{itemize}
        \item Distinguish the four classic definitions of AI.
        \item Explain why the \textbf{rational agent} view is central in modern AI.
        \item Specify a task using \textbf{PEAS} with concrete metrics.
        \item Classify environments (observable, Random, sequential, etc.).
        \item Choose an appropriate agent structure for a given problem.
    \end{itemize}
\end{frame}

% --- SECTION 1: INTRODUCTION ---
\section{Chapter 1: Defining the Field}

\begin{frame}{What is Artificial Intelligence?}
    AI is historically viewed through four distinct lenses:
    \vspace{0.4cm}
    \begin{columns}[T,onlytextwidth]
        \column{0.45\textwidth}
            \alert{Human-Centric}
            \begin{itemize}
                \item \textbf{Thinking Humanly:} Cognitive modeling.
                \item \textbf{Acting Humanly:} The Turing Test.
            \end{itemize}
        \column{0.45\textwidth}
            \alert{Rationality-Centric}
            \begin{itemize}
                \item \textbf{Thinking Rationally:} Laws of thought (Logic).
                \item \textbf{Acting Rationally:} The Rational Agent.
            \end{itemize}
    \end{columns}
    \vspace{0.5cm}
    \textit{This course focuses on the Rational Agent approach: acting to achieve the best outcome.}
\end{frame}

\begin{frame}{Four Views of AI: Strengths and Limits}
    \footnotesize
    \setlength{\tabcolsep}{4pt}
    \renewcommand{\arraystretch}{1.15}
    \begin{tabularx}{\textwidth}{>{\raggedright\arraybackslash}p{2.7cm} >{\raggedright\arraybackslash}X >{\raggedright\arraybackslash}X}
        \toprule
        \textbf{View} & \textbf{Main Idea} & \textbf{Typical Limitation} \\
        \midrule
        Thinking Humanly & Match human cognitive processes & Hard to validate internal thought processes \\
        Acting Humanly & Indistinguishable behavior from humans & Human-like behavior is not always optimal \\
        Thinking Rationally & Derive correct conclusions via logic & Real-world knowledge is incomplete/uncertain \\
        Acting Rationally & Select actions that maximize expected performance & Requires modeling uncertainty and trade-offs \\
        \bottomrule
    \end{tabularx}
\end{frame}

\begin{frame}{A Brief Historical Arc of AI}
    \begin{itemize}
        \item \textbf{1950s--1960s:} Symbolic AI and early optimism.
        \item \textbf{1970s--1980s:} Knowledge-based systems and expert systems.
        \item \textbf{1990s--2000s:} Probabilistic reasoning and machine learning become central.
        \item \textbf{2010s--Today:} Deep learning, foundation models, and hybrid systems.
    \end{itemize}
    \vspace{0.3cm}
    \textit{Book takeaway: progress accelerates when representation, learning, and computation improve together.}
\end{frame}

\begin{frame}{The Foundations of AI}
    AI is a "synthetic" discipline, drawing from:
    \begin{itemize}
        \item \textbf{Mathematics:} Probability (uncertainty) and Logic (certainty).
        \item \textbf{Economics:} Decision theory and utility maximization.
        \item \textbf{Neuroscience:} Hardware for intelligence (biological).
        \item \textbf{Control Theory:} Designing systems that maximize an objective function over time.
    \end{itemize}
\end{frame}

\begin{frame}{Reference Reading for This Part}
    \begin{itemize}
        \item \textbf{Russell \& Norvig (4th ed.):} Chapter 1 (Introduction), Chapter 2 (Intelligent Agents).
        \item \textbf{Poole \& Mackworth (3rd ed.):} Chapters on agents, environments, and problem formulation.
    \end{itemize}
    \vspace{0.3cm}
    Suggested strategy:
    \begin{enumerate}
        \item Read definitions first (AI views, rationality, PEAS).
        \item Then map each definition to one practical example.
        \item Finally, compare which assumptions break in real-world settings.
    \end{enumerate}
\end{frame}

% --- SECTION 2: INTELLIGENT AGENTS ---
\section{Chapter 2: The Agent Framework}

\begin{frame}{The Agent-Environment Interaction}
    An \textbf{Agent} is anything that:
    \begin{enumerate}
        \item Perceives its \textbf{Environment} through \textbf{Sensors}.
        \item Acts upon that environment through \textbf{Actuators}.
    \end{enumerate}
    \vspace{0.3cm}
    \begin{center}
        \fbox{Percepts $\rightarrow$ [ Agent Function ] $\rightarrow$ Actions}
    \end{center}
    \vspace{0.3cm}
    \textbf{Agent Program:} The concrete implementation of the mathematical Agent Function.
\end{frame}

\begin{frame}{Formal Agent Definition}
    Let:
    \begin{itemize}
        \item $\mathcal{P}$ be the set of possible percepts,
        \item $\mathcal{A}$ be the set of possible actions.
    \end{itemize}
    Then an agent function is:
    \[
        f: \mathcal{P}^{*} \rightarrow \mathcal{A}
    \]
    \vspace{0.2cm}
    Interpretation:
    \begin{itemize}
        \item Input: entire percept history.
        \item Output: next action.
    \end{itemize}
\end{frame}

\begin{frame}{Designing for Success: PEAS}
    Before writing code, we must specify the \textbf{PEAS} for the task:
    \begin{description}
        \item[Performance] The metric for success (e.g., safety, speed, profit).
        \item[Environment] The "world" (e.g., city streets, a chessboard).
        \item[Actuators] The tools for action (e.g., steering wheel, display).
        \item[Sensors] The inputs (e.g., cameras, keyboard, GPS).
    \end{description}
\end{frame}

\begin{frame}{PEAS Example: Autonomous Taxi}
    \small
    \begin{tabular}{p{2.2cm}p{8.1cm}}
        \toprule
        \textbf{Component} & \textbf{Example Specification} \\
        \midrule
        Performance & Safety, legality, trip time, passenger comfort, energy efficiency \\
        Environment & Urban/suburban roads, pedestrians, weather, traffic rules, other drivers \\
        Actuators & Steering, throttle, brake, indicators, display, audio prompts \\
        Sensors & Cameras, LiDAR, radar, GPS, IMU, wheel encoders, microphone \\
        \bottomrule
    \end{tabular}
\end{frame}

\begin{frame}{PEAS Example: Medical Diagnosis Assistant}
    \small
    \begin{tabular}{p{2.2cm}p{8.1cm}}
        \toprule
        \textbf{Component} & \textbf{Example Specification} \\
        \midrule
        Performance & Diagnostic accuracy, calibration, turnaround time, patient safety \\
        Environment & Hospital workflow, EHR systems, lab pipelines, changing clinical guidelines \\
        Actuators & Ranked differential list, recommendations, alerts, report generation \\
        Sensors & Symptoms, vitals, labs, imaging reports, history, clinician feedback \\
        \bottomrule
    \end{tabular}
\end{frame}

\begin{frame}{Task Environment Properties}
    The difficulty of an AI problem is defined by its environment:
    \begin{itemize}
        \item \textbf{Observable:} Fully vs. Partially (Can we see everything?).
        \item \textbf{Deterministic:} vs. Random (Is the next state predictable?).
        \item \textbf{Episodic:} vs. Sequential (Do current actions affect the future?).
        \item \textbf{Static:} vs. Dynamic (Does the world change while we think?).
        \item \textbf{Discrete:} vs. Continuous (Chess vs. Driving).
    \end{itemize}
\end{frame}

\begin{frame}{Environment Classification by Example}
    \footnotesize
    \setlength{\tabcolsep}{4pt}
    \renewcommand{\arraystretch}{1.12}
    \begin{tabularx}{\textwidth}{>{\raggedright\arraybackslash}p{2.5cm} *{4}{>{\raggedright\arraybackslash}X}}
        \toprule
        \textbf{Task} & \textbf{Observability} & \textbf{Dynamics} & \textbf{Outcome Model} & \textbf{State Type} \\
        \midrule
        Chess & Fully observable & Static (turn-based) & Mostly deterministic & Discrete \\
        Poker & Partially observable & Static (turn-based) & Random (hidden cards) & Discrete \\
        Robot vacuum & Partially observable & Dynamic & Random & Hybrid \\
        Highway driving & Partially observable & Dynamic & Random & Continuous \\
        \bottomrule
    \end{tabularx}
\end{frame}

\begin{frame}{The Hierarchy of Agent Structures}
    Modern AI agents are categorized by their complexity:
    \begin{itemize}
        \item \textbf{Simple Reflex:} Action based only on current percept.
        \item \textbf{Model-Based:} Maintains internal state to track the "unseen."
        \item \textbf{Goal-Based:} Acts to reach a specific destination or state.
        \item \textbf{Utility-Based:} Acts to maximize "happiness" (useful for trade-offs).
        \item \textbf{Learning Agents:} Can improve their performance over time.
    \end{itemize}
\end{frame}

\begin{frame}{When to Choose Each Agent Type}
    \begin{itemize}
        \item \textbf{Simple Reflex:} Use when environments are fully observable and rules are stable.
        \item \textbf{Model-Based:} Use when some state is hidden and must be inferred.
        \item \textbf{Goal-Based:} Use when planning to a target state is required.
        \item \textbf{Utility-Based:} Use when multiple goals conflict and trade-offs matter.
        \item \textbf{Learning Agent:} Use when environment is complex, non-stationary, or hard-coded rules fail.
    \end{itemize}
\end{frame}

\begin{frame}{Learning Agent Architecture (Conceptual)}
    A standard learning agent can be decomposed into:
    \begin{itemize}
        \item \textbf{Performance Element:} Chooses external actions.
        \item \textbf{Learning Element:} Improves the performance element from experience.
        \item \textbf{Critic:} Evaluates behavior against the performance measure.
        \item \textbf{Problem Generator:} Proposes exploratory actions to gather informative experience.
    \end{itemize}
\end{frame}

\begin{frame}{Key Takeaways}
    \begin{enumerate}
        \item Rationality is about \textbf{expected performance}, not perfect outcomes.
        \item PEAS is the first engineering step before implementation.
        \item Environment properties determine algorithmic difficulty.
        \item Agent architecture should match uncertainty, goals, and learning needs.
    \end{enumerate}
\end{frame}

\begin{frame}[standout]
    Questions?
\end{frame}

\end{document}